% GENERATED FILE. Edit canonical lessons, then run npm run generate:books.
% canonical-source-hash: fnv1a64:9d28341cb210601a
% canonical-lessons: JA-C363-kaikei, JA-C363-daikin, JA-C363-shouhizei, JA-C363-hangaku, JA-C363-zeitaku

\chapter{Paying the Bill}
\label{ch:ja-paying-the-bill}
\hlchaptermodality{363}

\begin{chapteropening}
\textbf{By the end of this chapter:} \emph{I can ask for the bill and talk about what things cost.}

Complete the last lesson of chapter 363: I can ask for the bill and talk about what things cost.
\end{chapteropening}

\section[\texorpdfstring{\ja{かいけい} (kaikei)}{かいけい (kaikei)}]{\ja{かいけい} (kaikei) — the bill; paying at the till}

\label{lesson:JA-C363-kaikei}



\begin{quote}
Before the new one: say the Japanese for prices in general, then the Japanese for a rental deposit.
\end{quote}

\subsection*{You'll want to know: \ja{かいけい}}
\textbf{\ja{かいけい}} — \emph{kaikei} — \textquotedblleft{}the bill; paying at the till\textquotedblright{}.

\begin{quote}
\textbf{\ja{しょくじが} \ja{おわったら、かいけいを} \ja{おねがいします。}} — \emph{shokuji ga owattara, kaikei o onegai shimasu} — \textquotedblleft{}When we have finished eating, the bill, please.\textquotedblright{}
\end{quote}

\begin{cousinweb}
\textbf{\ja{かいけい}} (\emph{kaikei}) is Sino-Japanese (kango): two Chinese characters read with their on'yomi, readings borrowed from Chinese. They are \emph{kai}, \textquotedblleft{}meeting\textquotedblright{} (Middle Chinese \emph{hwajH}) and \emph{kei}, \textquotedblleft{}plan\textquotedblright{} (\emph{kejH}) (KANJIDIC2; Middle Chinese in Baxter's transcription).
\end{cousinweb}

\begin{grammarlens}[title={What you've built}]
One more word for saying what happens if.
\end{grammarlens}

\subsection*{Your turn}
\begin{itemize}
\raggedright
  \item \emph{Say it:} \emph{kaikei}
  \item \emph{Say it:} \emph{kaikei}, once more
  \item \emph{Say it:} \emph{shokuji ga owattara, kaikei o onegai shimasu} — When we have finished eating, the bill, please.
  \item \emph{Recall:} say the Japanese for prices in general, then the Japanese for a rental deposit, then say \emph{kaikei} again
\end{itemize}

\begin{tcolorbox}[breakable,colback=teal!4,colframe=teal!35!black,title={Before you move on}]
What does \ja{かいけい} mean? (\textquotedblleft{}The bill; paying at the till\textquotedblright{}.) Say it once more.
\end{tcolorbox}
\section[\texorpdfstring{\ja{だいきん} (daikin)}{だいきん (daikin)}]{\ja{だいきん} (daikin) — the payment; the price of goods}

\label{lesson:JA-C363-daikin}



\begin{quote}
Before the new one: say the Japanese for a rental deposit, then the Japanese for the bill.
\end{quote}

\subsection*{You'll want to know: \ja{だいきん}}
\textbf{\ja{だいきん}} — \emph{daikin} — \textquotedblleft{}the payment; the price of goods\textquotedblright{}.

\begin{quote}
\textbf{\ja{しなものが} \ja{とどいたら、だいきんを} \ja{はらいます。}} — \emph{shinamono ga todoitara, daikin o haraimasu} — \textquotedblleft{}When the goods arrive, I will pay for them.\textquotedblright{}
\end{quote}

\begin{cousinweb}
\textbf{\ja{だいきん}} (\emph{daikin}) is Sino-Japanese (kango): two Chinese characters read with their on'yomi, readings borrowed from Chinese. They are \emph{dai}, \textquotedblleft{}fee\textquotedblright{} (Middle Chinese \emph{dojH}) and \emph{kin}, \textquotedblleft{}gold\textquotedblright{} (\emph{kim}) (KANJIDIC2; Middle Chinese in Baxter's transcription).
\end{cousinweb}

\begin{grammarlens}[title={What you've built}]
One more word for saying what happens if.
\end{grammarlens}

\subsection*{Your turn}
\begin{itemize}
\raggedright
  \item \emph{Say it:} \emph{daikin}
  \item \emph{Say it:} \emph{daikin}, once more
  \item \emph{Say it:} \emph{shinamono ga todoitara, daikin o haraimasu} — When the goods arrive, I will pay for them.
  \item \emph{Recall:} say the Japanese for a rental deposit, then the Japanese for the bill, then say \emph{daikin} again
\end{itemize}

\begin{tcolorbox}[breakable,colback=teal!4,colframe=teal!35!black,title={Before you move on}]
What does \ja{だいきん} mean? (\textquotedblleft{}The payment; the price of goods\textquotedblright{}.) Say it once more.
\end{tcolorbox}
\section[\texorpdfstring{\ja{しょうひぜい} (shōhizei)}{しょうひぜい (shōhizei)}]{\ja{しょうひぜい} (shōhizei) — consumption tax}

\label{lesson:JA-C363-shouhizei}



\begin{quote}
Before the new one: say the Japanese for the bill, then the Japanese for the payment.
\end{quote}

\subsection*{You'll want to know: \ja{しょうひぜい}}
\textbf{\ja{しょうひぜい}} — \emph{shōhizei} — \textquotedblleft{}consumption tax\textquotedblright{}.

\begin{quote}
\textbf{\ja{しょうひぜいが} \ja{あがったら、かいものを} \ja{ひかえるでしょう。}} — \emph{shōhizei ga agattara, kaimono o hikaeru deshō} — \textquotedblleft{}If consumption tax goes up, people will probably cut down on shopping.\textquotedblright{}
\end{quote}

\begin{cousinweb}
\textbf{\ja{しょうひぜい}} (\emph{shōhizei}) is Sino-Japanese (kango): three Chinese characters read with their on'yomi, readings borrowed from Chinese. They are \emph{shō}, \textquotedblleft{}extinguish\textquotedblright{} (Middle Chinese \emph{sjew}), \emph{hi}, \textquotedblleft{}expense\textquotedblright{} (\emph{phj+jH}) and \emph{zei}, \textquotedblleft{}tax\textquotedblright{} (\emph{sywejH}) (KANJIDIC2; Middle Chinese in Baxter's transcription).
\end{cousinweb}

\begin{grammarlens}[title={What you've built}]
One more word for saying what happens if.
\end{grammarlens}

\subsection*{Your turn}
\begin{itemize}
\raggedright
  \item \emph{Say it:} \emph{shōhizei}
  \item \emph{Say it:} \emph{shōhizei}, once more
  \item \emph{Say it:} \emph{shōhizei ga agattara, kaimono o hikaeru deshō} — If consumption tax goes up, people will probably cut down on shopping.
  \item \emph{Recall:} say the Japanese for the bill, then the Japanese for the payment, then say \emph{shōhizei} again
\end{itemize}

\begin{tcolorbox}[breakable,colback=teal!4,colframe=teal!35!black,title={Before you move on}]
What does \ja{しょうひぜい} mean? (\textquotedblleft{}Consumption tax\textquotedblright{}.) Say it once more.
\end{tcolorbox}
\section[\texorpdfstring{\ja{はんがく} (hangaku)}{はんがく (hangaku)}]{\ja{はんがく} (hangaku) — half price}

\label{lesson:JA-C363-hangaku}



\begin{quote}
Before the new one: say the Japanese for the payment, then the Japanese for consumption tax.
\end{quote}

\subsection*{You'll want to know: \ja{はんがく}}
\textbf{\ja{はんがく}} — \emph{hangaku} — \textquotedblleft{}half price\textquotedblright{}.

\begin{quote}
\textbf{\ja{ばんに} \ja{なったら、おべんとうが} \ja{はんがくに} \ja{なります。}} — \emph{ban ni nattara, obentō ga hangaku ni narimasu} — \textquotedblleft{}In the evening, the boxed lunches go half price.\textquotedblright{}
\end{quote}

\begin{cousinweb}
\textbf{\ja{はんがく}} (\emph{hangaku}) is Sino-Japanese (kango): two Chinese characters read with their on'yomi, readings borrowed from Chinese. They are \emph{han}, \textquotedblleft{}half\textquotedblright{} (Middle Chinese \emph{panH}) and \emph{gaku}, \textquotedblleft{}sum\textquotedblright{} (\emph{ngaek}) (KANJIDIC2; Middle Chinese in Baxter's transcription).
\end{cousinweb}

\begin{grammarlens}[title={What you've built}]
One more word for saying what happens if.
\end{grammarlens}

\subsection*{Your turn}
\begin{itemize}
\raggedright
  \item \emph{Say it:} \emph{hangaku}
  \item \emph{Say it:} \emph{hangaku}, once more
  \item \emph{Say it:} \emph{ban ni nattara, obentō ga hangaku ni narimasu} — In the evening, the boxed lunches go half price.
  \item \emph{Recall:} say the Japanese for the payment, then the Japanese for consumption tax, then say \emph{hangaku} again
\end{itemize}

\begin{tcolorbox}[breakable,colback=teal!4,colframe=teal!35!black,title={Before you move on}]
What does \ja{はんがく} mean? (\textquotedblleft{}Half price\textquotedblright{}.) Say it once more.
\end{tcolorbox}
\section[\texorpdfstring{\ja{ぜいたく} (zeitaku)}{ぜいたく (zeitaku)}]{\ja{ぜいたく} (zeitaku) — luxury; extravagant}

\label{lesson:JA-C363-zeitaku}



\begin{quote}
Before the new one: say the Japanese for consumption tax, then the Japanese for half price.
\end{quote}

\subsection*{You'll want to know: \ja{ぜいたく}}
\textbf{\ja{ぜいたく}} — \emph{zeitaku} — \textquotedblleft{}luxury; extravagant\textquotedblright{}.

\begin{quote}
\textbf{\ja{おかねが} \ja{あったら、ぜいたくを} \ja{したいです。}} — \emph{okane ga attara, zeitaku o shitai desu} — \textquotedblleft{}If I had the money, I would like to live it up.\textquotedblright{}
\end{quote}

A na-adjective: \ja{な} before a noun.

\begin{cousinweb}
\textbf{\ja{ぜいたく}} (\emph{zeitaku}) is Sino-Japanese (kango): two Chinese characters read with their on'yomi, readings borrowed from Chinese. They are \emph{zei}, \textquotedblleft{}luxury\textquotedblright{} (Middle Chinese \emph{tsywejH}) and \emph{taku}, \textquotedblleft{}swamp\textquotedblright{} (KANJIDIC2; Middle Chinese in Baxter's transcription).
\end{cousinweb}

\begin{grammarlens}[title={What you've built}]
One more word for saying what happens if.
\end{grammarlens}

\subsection*{Your turn}
\begin{itemize}
\raggedright
  \item \emph{Say it:} \emph{zeitaku}
  \item \emph{Say it:} \emph{zeitaku}, once more
  \item \emph{Say it:} \emph{okane ga attara, zeitaku o shitai desu} — If I had the money, I would like to live it up.
  \item \emph{Recall:} say the Japanese for consumption tax, then the Japanese for half price, then say \emph{zeitaku} again
\end{itemize}

\begin{tcolorbox}[breakable,colback=teal!4,colframe=teal!35!black,title={Before you move on}]
What does \ja{ぜいたく} mean? (\textquotedblleft{}Luxury; extravagant\textquotedblright{}.) Say it once more.
\end{tcolorbox}
